\section{A positive integral for the homogeneous moment}\label{sec:integral}

The arch formula \eqref{eq:01-arch-formula} is exact, but its Pfaffian
quotient does not directly show the size of $m_N$. We convert the
coalescing-parameter problem into a positive integral. This will establish
the limiting quotient without assuming that its denominator remains
nonzero, and make the moment accessible to comparison estimates.
For generic $X$, put $b=\cos w$ and
\begin{equation}\label{eq:01-Q-definition}
 \begin{split}
 Q_X(b)&=(b-C)g_X(b)\prod_{i=1}^N
  (b-a_i)(\cos(w-2\lambda)-a_i)(\cos(w+2\lambda)-a_i),\\
 \Phi_X(w)&=\sin w\,Q_X(\cos w).
 \end{split}
\end{equation}
The paired shifted factors have product
$b^2-2da_i b+a_i^2-d^2$, so $Q_X$ is a polynomial of degree $4N+1$.

\begin{lemma}[Missing harmonics]\label{lem:missing-harmonics}
Every sine harmonic of $\Phi_X$ whose frequency is divisible by four
has coefficient zero. Equivalently, writing
$h=b^2(1-b^2)=(1-\cos4w)/8$, there are polynomials such that
\begin{equation}\label{eq:01-projection}
 Q_X(b)=\mathcal P(h)+b\mathcal J(h)+b^2\mathcal E(h),
 \quad \deg\mathcal P,\deg\mathcal J\le N,
 \quad \deg\mathcal E\le N-1.
\end{equation}
\end{lemma}

\begin{proof}
Fix $i$ and write $f_j=\cos(v+j\lambda)$ with $f_0=a_i$.
The product in \eqref{eq:01-Q-definition} gives
$\Phi_X(v\pm2\lambda)=0$. The deletion quotient also gives
$\Phi_X(v+6\lambda)+\Phi_X(v-6\lambda)=0$. To verify its sign and
factors, cancel the contributions of the other $a_j$ in the ratio:
on both sides these contributions involve exactly
$f_6,f_{-6},f_4,f_{-4},-f_0$. The remaining expressions are
\[
 R_\pm=\sin(v\pm6\lambda)(f_{\pm6}-C)
       (f_0+f_{\pm6}+\rho)(f_{\pm6}-f_0)(f_{\pm4}-f_0).
\]
Product-to-sum and $\rho=2cd$ give
\[
 R_\pm=\cos(2v)(f_{\pm6}-C)
             \sin(v\pm3\lambda)(\cos(v\pm3\lambda)+d).
\]
The last two expressions after removing $\cos(2v)$ are trigonometric
polynomials of degree three. They have the same six zeros
$v=\pm3\lambda,\pm5\lambda,\pm7\lambda$ modulo $2\pi$, and
opposite values at zero; hence they are negatives of each other.
This proves the claimed cancellation as a polynomial identity,
including where canceled factors vanish.

Consider the sum of the four shifts of $\Phi_X$ by
$\pm2\lambda,\pm6\lambda$. On frequency $k$ its multiplier is
$4\cos(4k\lambda)\cos(2k\lambda)$, nonzero precisely when four divides
$k$. Its maximum frequency is $4N+2$. Dividing this shifted sum by
$\sin4w$ therefore gives a polynomial in $\cos4w$ of degree at most
$N-1$. The preceding identities give $N$ distinct roots, for generic
$a_i$, so it is zero. This proves the harmonic assertion.

Every polynomial in $b$ has a unique expression
$P_0(h)+bP_1(h)+b^2P_2(h)+b^3P_3(h)$, by repeatedly using
$b^4=b^2-h$. The first three terms, after multiplication by $\sin w$,
have no frequency divisible by four. For the last term use
\[
 \sin w\,b^3=\frac14\sin2w+\frac18\sin4w.
\]
Its projection onto multiples of four is
$\frac18\sin4w\,P_3((1-\cos4w)/8)$, so it vanishes only if $P_3=0$.
The degree bounds follow from $\deg Q_X=4N+1$.
\end{proof}

At the homogeneous parameters $a_i=C$, the paired factors in
\eqref{eq:01-Q-definition} become $(b-C)(b-c)$. Thus the desired
limiting monic polynomial is characterized by
\begin{equation}\label{eq:01-homogeneous-projection}
 Q(b)=(b-C)^{2N+1}(b-c)^N g(b)
     =\mathcal P(h)+b\mathcal J(h)+b^2\mathcal E(h),
 \qquad \deg g=N.
\end{equation}

\begin{lemma}[Nonsingularity and integral representation]
\label{lem:positive-integral}
The monic problem \eqref{eq:01-homogeneous-projection} has a unique
solution, which is the limit of $g_X$ as all $a_i\to C$.
Let
\begin{equation}\label{eq:01-omega}
 \omega(x)=\frac12\sqrt{\sqrt{\frac{1+x}{2x}}-1},\qquad 0<x<1,
\end{equation}
and let $(y_1,\ldots,y_N)$ have density proportional to
\begin{equation}\label{eq:01-positive-law}
 \prod_{i=1}^N\omega(y_i)
 \prod_{i<j}\frac{(y_j-y_i)^2}{y_i+y_j},
 \qquad 0<y_1<\cdots<y_N<1.
\end{equation}
Then
\begin{equation}\label{eq:01-moment-integral}
 m_N=\frac\eta\pi\int_0^1
  \left(1-\E\prod_{i=1}^N\frac{y_i-x}{y_i+x}\right)
  \frac{\omega(x)}x\,dx.
\end{equation}
\end{lemma}

\begin{proof}
We first turn the zero-multiplicity conditions into Cauchy-transform
moment equations. A positive determinant will make those equations
uniquely solvable and control the coalescing limit. We then identify the
resulting coefficient with the expectation in
\eqref{eq:01-moment-integral}.
The proposed probability law is well defined: $\omega(x)=O(x^{-1/4})$
at zero and each pair factor is at most $y_j\le1$, so its normalizing
integral is finite; it is positive because the density is positive
inside the ordered region.

\smallskip\noindent\emph{Cauchy moment equations.}
The two prescribed zeros satisfy $h(C)=h(c)=1/8$. To place them at
infinity, set $z=(8h-1)^{-1}$, so that the inverse relation is
\[
 b^2(1-b^2)=\frac{1+1/z}{8}.
\]
We use the two branches near infinity tending to $C$ and $c$, respectively.
Their differences from these limits have order $z^{-1}$, since
$h'(C)$ and $h'(c)$ are nonzero. For each polynomial in
\eqref{eq:01-homogeneous-projection}, multiply by $z^N$ after the
substitution $h=(1+1/z)/8$. Denote the resulting polynomials by
$\widetilde{\mathcal P}$, $\widetilde{\mathcal J}$, and
$\widetilde{\mathcal E}$; the last satisfies
$\widetilde{\mathcal E}(0)=0$ because $\deg\mathcal E\le N-1$.
The branch tending to $C$ and the sum of the two branches are, respectively,
\[
 f(z)=\sqrt{\frac{1+\sqrt{(1-1/z)/2}}2},\qquad
 l(z)=\sqrt{1+\sqrt{(1+1/z)/2}}.
\]
These formulas continue $f$ analytically off $[0,1]$ and $l$ off
$[-1,0]$. On the $C$ branch, the zero of order $2N+1$ in
\eqref{eq:01-homogeneous-projection} becomes $O(z^{-2N-1})$; on the
$c$ branch, the zero of order $N$ becomes $O(z^{-N})$.
After multiplication by $z^N$ this gives
\begin{align}
 \mathcal F(z)&=\widetilde{\mathcal P}(z)
        +f(z)\widetilde{\mathcal J}(z)
        +f(z)^2\widetilde{\mathcal E}(z)=O(z^{-N-1}),
        \label{eq:01-F-decay}\\
 \mathcal D(z)&=\widetilde{\mathcal J}(z)
                 +l(z)\widetilde{\mathcal E}(z)=O(1).
        \label{eq:01-D-decay}
\end{align}
For the second statement, evaluate the same polynomial expression
on the branch tending to $c$, where multiplication by $z^N$ makes
it $O(1)$, and subtract the expression on the $C$ branch; the
quotient by their nonzero difference is $\mathcal D$.

With subscripts $+$ and $-$ denoting boundary values from the upper
and lower half-planes, the jumps across the cuts are
\[
 f_+(x)-f_-(x)=2i\omega(x),\qquad
 l_+(-x)-l_-(-x)=-2\sqrt2\,i\omega(x),\qquad
 f_+(x)+f_-(x)=l(x).
\]
At the outer endpoints these functions are bounded. At zero,
$\mathcal F=O(|z|^{-1/4})$ and
$\mathcal D(z)\to D_0:=\widetilde{\mathcal J}(0)$;
the latter uses $\widetilde{\mathcal E}(0)=0$. Cauchy's formula,
with small endpoint circles shrinking to zero, yields
\begin{align}
 \mathcal F(z)&=\frac1\pi\int_0^1
                   \frac{\mathcal D(x)\omega(x)}{x-z}\,dx,
        \label{eq:01-F-Cauchy}\\
 \int_0^1x^j\mathcal D(x)\omega(x)\,dx&=0,
                  \qquad 0\le j<N,
        \label{eq:01-moments}\\
 \mathcal D(z)&=D_0-\frac{\sqrt2}{\pi}z
       \int_0^1\frac{\widetilde{\mathcal E}(-y)/y}{z+y}
                    \omega(y)\,dy.
        \label{eq:01-D-Cauchy}
\end{align}
The moment equations are the first $N$ coefficients of the expansion
of \eqref{eq:01-F-Cauchy} at infinity. To obtain
\eqref{eq:01-D-Cauchy}, subtract the bounded-at-infinity representation
at zero. Its integrand is integrable there because
$\widetilde{\mathcal E}(-y)/y$ is a polynomial.

To control the endpoint at zero, use $\omega(x)=O(x^{-1/4})$ and
$\widetilde{\mathcal E}(z)=O(z)$ near zero. The increment in
\eqref{eq:01-D-Cauchy}, split at $y=x$, is bounded by
\[
 C\left(\int_0^x y^{-1/4}\,dy
       +x\int_x^1y^{-5/4}\,dy\right)=O(x^{3/4}).
\]
This estimate will also justify the singular moment integral below.

\smallskip\noindent\emph{Nonsingularity and coalescence.}
The polynomial $\widetilde{\mathcal E}(-y)/y$ has degree at most
$N-1$, and the moment equations determine it uniquely from $D_0$.
Indeed the corresponding matrix, apart from nonzero scalar factors,
is
\begin{equation}\label{eq:01-moment-matrix}
 M_{jk}=\int_0^1\int_0^1
           \frac{x^{j+1}y^k}{x+y}\omega(x)\omega(y)\,dx\,dy,
           \qquad 0\le j,k<N.
\end{equation}
Its determinant is strictly positive. This is the Cauchy-bimoment
positivity argument used in
\cite[Theorem~2.1]{BertolaGekhtmanSzmigielski2010}; we give the finite
determinant calculation for the present weights. Expand the determinant and
antisymmetrize successively in the $x$ and $y$ variables. On the
ordered integration regions the integrand becomes the product of
two positive Vandermonde determinants, $\prod_i x_i$, positive
weights, and the Cauchy determinant
\[
 \det\left[\frac1{x_i+y_j}\right]
 =\frac{\prod_{i<j}(x_j-x_i)(y_j-y_i)}{\prod_{i,j}(x_i+y_j)}>0.
\]
For completeness this identity follows by multiplying through by the
denominator: the numerator is alternating in each family and has
exactly the degree of the two Vandermondes. The remaining constant
is one, recursively evaluating at $x_N=-y_N$.

Consider now a solution of the homogeneous linear problem for $g$,
so $\deg g\le N-1$. The coefficient of $b^{4N+1}$ in $Q$ is zero.
Since
\[
 [b^{4N+1}]Q=(-8)^N D_0,
\]
we obtain $D_0=0$. Positivity of \eqref{eq:01-moment-matrix} forces
$\widetilde{\mathcal E}=0$; then
\eqref{eq:01-D-Cauchy} gives $\widetilde{\mathcal J}=0$, and
\eqref{eq:01-F-decay} gives $\widetilde{\mathcal P}=0$. Thus $g=0$.
For the monic problem, write $g(b)=b^N+\sum_{r=0}^{N-1}\gamma_r b^r$.
The excluded $b^3P_3(h)$ part in the proof of
Lemma~\ref{lem:missing-harmonics} has $N$ coefficients. Requiring them
to vanish therefore gives an $N\times N$ linear system for
$(\gamma_0,\ldots,\gamma_{N-1})$, and the established injectivity
proves that this system is invertible at the homogeneous parameters.
This proves existence and uniqueness of $g$.
Apply the same coefficient conditions to \eqref{eq:01-Q-definition},
with $g_X$ replaced by a monic unknown polynomial. Their matrix and right side depend
polynomially on the $a_i$, through the displayed zero factors, and
converge to those of the homogeneous system. This coefficient matrix
is consequently invertible for all $a_i$ sufficiently close to $C$.
The generic $g_X$ satisfies that system by
Lemma~\ref{lem:missing-harmonics}, so it converges to $g$.
The moment itself is regular at the homogeneous point by
Lemma~\ref{lem:transfer-decay}.

\smallskip\noindent\emph{The moment coefficient and its expectation.}
For this monic solution $D_0=(-8)^{-N}$. Put
$G(z)=\mathcal D(z)/D_0$.
The leading coefficient computation also gives
$[b^{4N}]Q=(-8)^N\widetilde{\mathcal P}(0)$. Comparing this with the
product in \eqref{eq:01-homogeneous-projection} gives the coefficient
needed by the arch formula:
\[
 [b^{N-1}]g(b)=\frac{\widetilde{\mathcal P}(0)}{D_0}
                    +(2N+1)C+Nc.
\]
Substitution into \eqref{eq:physical-arch-coefficient}, using
$C=\eta(2C+c)$, yields the first equality below:
\begin{equation}\label{eq:01-moment-D}
 m_N=-\eta\left(C+\frac{\widetilde{\mathcal P}(0)}{D_0}\right)
     =\frac\eta\pi\int_0^1(1-G(x))\frac{\omega(x)}x\,dx.
\end{equation}
For the second equality, use
\[
 f(z)=C+\frac1\pi\int_0^1\frac{\omega(x)}{x-z}\,dx
\]
and subtract $D_0f$ from \eqref{eq:01-F-Cauchy}, then let $z\to0-$.
Equation \eqref{eq:01-D-Cauchy} gives
$\mathcal D(x)-D_0=O(x^{3/4})$, so the integrand after subtraction is $O(x^{-1/2})$ at zero.
This ensures integrability and justifies the limit. On the left, terms involving
$f(z)(\widetilde{\mathcal J}(z)-D_0)$ and
$f(z)^2\widetilde{\mathcal E}(z)$ tend to zero.

The moment equations now have a unique solution. It remains to
exhibit their probabilistic solution, thereby identifying $G$ with the
expectation in \eqref{eq:01-moment-integral}. Let $\mathcal W_n$ denote
the product in \eqref{eq:01-positive-law} with $n$ coordinates, viewed
on $(0,1)^n$ without the ordering constraint. We use its normalized
symmetric law for $n=N$. Partial fractions give
\[
 \prod_{i=1}^N\frac{y_i-x}{y_i+x}
 =1-x\sum_{i=1}^N\frac{\alpha_i}{x+y_i},
 \qquad
 \alpha_i=2\prod_{j\ne i}\frac{y_j+y_i}{y_j-y_i}.
\]
After multiplication by the density, the factors in $\alpha_i$
cancel one Vandermonde and all corresponding denominator factors.
Each canceled term is integrable: its remaining polynomial factors
and pair kernels are bounded on $(0,1)^N$, as is $x/(x+y_i)$, and
the one-variable weights are integrable.
Integrating over the other variables therefore gives a representation
$1-x\int_0^1P(y)\omega(y)/(x+y)\,dy$ with $\deg P\le N-1$.
Its moments against $x^k\omega(x)$ vanish for $k<N$. To see this,
adjoin $x$ as an $(N+1)$st variable and use the density identity
\[
 x^k\omega(x)\mathcal W_N(y_1,\ldots,y_N)
       \prod_{i=1}^N\frac{y_i-x}{y_i+x}
 =\frac{x^k\mathcal W_{N+1}(x,y_1,\ldots,y_N)}
          {\prod_{i=1}^N(y_i-x)}.
\]
The left side is absolutely integrable, since the product has modulus
at most one. Symmetrizing the right side in all $N+1$ variables gives
zero by the elementary interpolation identity
\[
 \sum_{r=0}^N\frac{z_r^k}{\prod_{j\ne r}(z_j-z_r)}=0,
 \qquad 0\le k<N.
\]
Thus the expectation has the same polynomial Cauchy representation
and vanishing moments as $G$. The uniqueness argument using
\eqref{eq:01-moment-matrix} now identifies this expectation with $G$.
Substitute into \eqref{eq:01-moment-D}.
\end{proof}
